\documentclass[12pt]{article}

\setlength{\parindent}{0pt}
\setcounter{secnumdepth}{3}
\usepackage[margin=1in]{geometry}

% packages
\usepackage{amsmath, amsfonts, amsthm, amssymb}
\usepackage{graphicx, float}
\usepackage{mathtools}
\usepackage{titlesec}
\usepackage{interval}
\usepackage{hyperref}
\usepackage{siunitx}
\usepackage{titling}
\usepackage{vwcol}
\usepackage{setspace}
\usepackage{empheq}
\usepackage{cancel}
\usepackage{esdiff}
\usepackage{multicol}
\usepackage{mdframed}
\usepackage{esdiff}
\usepackage{tikzsymbols}
\usepackage{multicol}
\usepackage{tikz}
\usepackage{varwidth}
\usepackage{pgfplots}
\usepackage{fancyhdr}
\usepackage{enumitem}
\usepackage{tcolorbox}

\intervalconfig {
	soft open fences
}

% header
\pagestyle{fancy}
\fancyhf{}
\lhead{Ethan Fan}
\rhead{AP Calculus BC}
\cfoot{\thepage}

% section format
\titleformat{\section}
  {\large \bfseries}
  {}
  {0em}
  {}
\titlespacing*{\section}{0pt}{0.5em}{0.3em}

\titleformat{\subsection}[runin]
  {\bfseries}
  {}
  {0em}
  {}
\titlespacing*{\subsection}{0pt}{0pt}{0em}

\titleformat{\subsubsection}[runin]
  {\itshape}
  {(\roman{subsubsection})}
  {0em}
  {}

% commands
\newcommand{\alignedintertext}[1]{%
  \noalign{%
    \vskip\belowdisplayshortskip
    \vtop{\hsize=\linewidth#1\par
    \expandafter}%
    \expandafter\prevdepth\the\prevdepth
  }%
}

\newcommand*{\paren}[1]{\ensuremath\left(#1\right)}
\newcommand*{\limit}[2][x]{\ensuremath{\displaystyle\lim_{#1\to#2}}}
\newcommand*{\Limit}[3][x]{\ensuremath{\displaystyle\lim_{#1\to#2}\left[#3\right]}}
\newcommand*{\deriv}[1][x]{\ensuremath{\dfrac{\mathrm{d}}{\mathrm{d}#1}}}
\newcommand*{\Deriv}[2][x]{\ensuremath{\dfrac{\mathrm{d}}{\mathrm{d}#1}\left[#2\right]}}
\newcommand{\ddx}{\dfrac{d}{dx}}
\newcommand{\dydx}{\dfrac{dy}{dx}}
\newcommand{\dydt}{\dfrac{dy}{dt}}
\newcommand{\dxdt}{\dfrac{dx}{dt}}
\newcommand*{\iinteg}[2][x]{\ensuremath{\displaystyle\int #2\;\mathrm{d}#1}}
\newcommand*{\dinteg}[4][x]{\ensuremath{\displaystyle\int_{#2}^{#3}#4\;\mathrm{d}#1}}
\newcommand*{\abs}[1]{\ensuremath{\left|#1\right|}}
\newcommand*{\eps}{\varepsilon}
\newcommand*{\real}{\ensuremath\mathbb{R}}
\newcommand*{\integer}{\ensuremath\mathbb{Z}}
\newcommand*{\posint}{\ensuremath\mathbb{Z^+}}
\newcommand*{\cbrt}[1]{\ensuremath{\sqrt[3]{#1}}}
\newcommand*{\lheqzero}{\ensuremath{\underset{\text{L'H}}{\overset{\left[\frac00\right]}{=}}}}
\newcommand*{\lheqinfty}{\ensuremath{\underset{\text{L'H}}{\overset{\left[\frac{\infty}{\infty}\right]}{=}}}}
\newcommand{\tor}{\text{ or }}
\newcommand{\floor}[1]{\lfloor #1 \rfloor}

\newcommand{\vem}{\vspace{0.5em}}
\newcommand{\example}[1]{\section{Example #1}}
\newcommand{\problem}[1]{\section{Problem #1}}
\newcommand{\subproblem}[1]{\subsection{(#1)} \,}

\DeclareMathOperator{\dne}{DNE}
\DeclareMathOperator{\sgn}{sgn}

\newcommand{\definition}[1]{\begin{tcolorbox}[colback=red!5!white,colframe=red!75!black,parbox=false] #1 \end{tcolorbox}}

\newcommand{\theorem}[2]{\begin{tcolorbox}[title={#1},colback=blue!5!white,colframe=blue!75!black,parbox=false] #2 \end{tcolorbox}}

\allowdisplaybreaks
% \postdisplaypenalty=100000

% document
\begin{document}

\vspace*{0cm}

% opening
\begin{center}
    {\Large \bfseries Problem Set 13} \\[0.5cm]
    {\large Integration by Parts I} \\[0.35cm]
    {\normalsize October 2, 2026}
\end{center}

% problems

\problem{2}

\subproblem{a}
\begin{gather*}
    \int xe^x\, dx = 
    \begin{bmatrix} e^x & x \,dx\\
    e^x \, dx& \frac{1}{2}x^2
    \end{bmatrix} = \frac{e^x}{2}x^2 - \int \frac{e^x}{2}x^2\, dx
\end{gather*}
The integral is harder to solve. \vem

\subproblem{b}
\begin{gather*}
    \int x^3 e^{x^2}\, dx =
    \begin{bmatrix}
        x^2 & xe^{x^2} \, dx\\
        2x\, dx & \frac{1}{2}e^{x^2}
    \end{bmatrix} 
    \frac{1}{2}x^2e^{x^2} - \int xe^{x^2} \, dx\\
    \int xe^{x^2} \, dx = \frac{1}{2}e^{x^2}\\
    \int x^3 e^{x^2}\, dx = \boxed{\frac{1}{2}(x^2-1)e^{x^2}+C}
\end{gather*}

\subproblem{c}
\begin{gather*}
    \int u\, dv = u(v+C) - \int v+C\, du=uv + uC - \int v \, du-uC = uv-\int v\, du\\
    \int \ln (x+1)\, dx= 
    \begin{bmatrix}
        \ln(x+1) & dx\\
        \frac{1}{x+1}\, dx & x+1
    \end{bmatrix}
    (x+1)\ln (x+1) - \int 1\, dx \\= \boxed{ (x+1)\ln (x+1)-x+C}
\end{gather*}

\subproblem{d}
The integrals on each side could hold different constants of integration, thereby differing the numerical value of each side of the equation by one.

\problem{3}

\subproblem{a}
The two areas under the curve equal the area formed by the large rectangle subtracting the small rectangle. \vem

\subproblem{b}
\begin{gather*}
    u=x=\sin y\\
    v=y=\arcsin x\\
    \int_0^1 \arcsin x\, dx+\int_0^{\pi /2} \sin y \, dy = \frac{\pi}{2}\\
    \int_0^1 \arcsin x\, dx = \frac{\pi }{2} - \int_0^{\pi /2} \sin y \, dy = \boxed{\frac{\pi}{2}-1}
\end{gather*}

\subproblem{c}
\begin{gather*}
    \int_1^e \ln x \, dx = e -\int_0^1 e^x\, dx = \boxed{1}
\end{gather*}

\problem{4}
\begin{gather*}
    f_{ave}=\frac{\int_1^3x^2\ln x\, dx}{2}\\
    \int_1^3x^2\ln x\, dx = 
    \begin{bmatrix}
        \ln x & x^2\, dx\\
        \frac{1}{x}\, dx & \frac{1}{3}x^3\,
    \end{bmatrix}
    \frac{1}{3}x^3\ln x\Big|^3_1 - \int_1^3 \frac{1}{3}x^2 \, dx\\
    = 9 \ln 3 + \frac{1}{9} - 3  = 9 \ln 3 - \frac{26}{9}\\
    f_{ave} = \frac{1}{2}(9 \ln 3 - \frac{26}{3}) = \boxed{\frac{9 \ln 3}{2}- \frac{13}{9}}
\end{gather*}

\problem{5}

\subproblem{a}
\begin{gather*}
    \int x^2 \sin x\, dx = \boxed{- x^2 \cos x + 2x \sin x + 2 \cos x +C}
\end{gather*}

\subproblem{b}
\begin{gather*}
    \int  e^x \sin x\, dx = e^x \sin x - e^x \cos x - \int  e^ x \sin x \, dx\\
    \int  e^x \sin x\, dx = \boxed{\frac{e^x (\sin x -\cos x )}{2}+C}
\end{gather*}

\subproblem{c}
\begin{gather*}
    \int x^3 e^{2x} \, dx = x^3 \cdot \frac{1}{2}e^{2x} - 3x^2 \cdot \frac{1}{4}e^{2x} + 6x \cdot \frac{1}{8}e^{2x} - 6 \cdot \frac{1}{16}e^{2x}\\
    = \boxed{\frac{e^{2x}}{8}(4x^3-6x^2+6x-3)+C}
\end{gather*}

\subproblem{d}
The table in (c) stops, as differentiating a polynomial will reach 0 eventually, while exponential and trig functions perpetuate. The method of (b) utilizes the fact that the original integral will appear on the other side after a finite number of integrations by parts. 

\problem{6}

\subproblem{c}
\begin{gather*}
    \int f^{-1}(x)\, dx =  
    \begin{bmatrix}
        f^{-1}(x) & dx \\
        (f^{-1})'(x)\, dx & x 
    \end{bmatrix}
    xf^{-1}(x) - \int x(f^{-1})'(x)\, dx\\
    \int x(f^{-1})'(x)\, dx = \int f(y)\, dy = F(y)=F(f^{-1}(x))\\
    \int f^{-1}(x)\, dx =  \boxed{xf^{-1}(x) - F(f^{-1}(x)) +C}
\end{gather*}

\subproblem{d}
\begin{gather*}
    f(x)=e^x \implies f^{-1}(x)=\ln x\\
    \int \ln x \, dx = x \ln x - e^{\ln x } +C= \boxed{x\ln x -x+C}
\end{gather*}

\problem{7}

\subproblem{b}
\begin{gather*}
    \int \sin x \cos 4x \, dx=
    \begin{bmatrix}
        \sin x & \cos 4 x\, dx\\
        \cos x \, dx& \frac{1}{4}\sin 4x
    \end{bmatrix}
    \frac{1}{4}\sin x \sin 4x-\frac{1}{4}\int \cos x \sin 4 x\, dx\\
    =\begin{bmatrix}
        \cos x & \sin 4 x\, dx\\
        -\sin x \, dx & -\frac{1}{4}\cos 4x
    \end{bmatrix}
    \frac{1}{4}\sin x \sin 4x + \frac{1}{16}\cos x \cos 4x + \frac{1}{16}\int \sin x \cos 4x\, dx\\
    \frac{15}{16}\int \sin x \cos 4x\, dx = \frac{1}{4}\sin x \sin 4x + \frac{1}{16}\cos x \cos 4x\\
    \int \sin x \cos 4x\, dx  = \boxed{\frac{4}{15}\sin x \sin 4x + \frac{1}{15}\cos x \cos 4x+C}
\end{gather*}
If the question is different, the $u$-sub may be applied. \vem

\subproblem{c}
\begin{gather*}
    \int \frac{xe^{2x}}{(2x+1)^2} \, dx =
    \begin{bmatrix}
        xe^{2x} & (2x+1)^{-2 } \, dx\\
        e^{2x}(2x+1)\, dx & - \frac{1}{2(2x+1)}
    \end{bmatrix}
    -\frac{xe^{2x}}{4x+2} + \int \frac{e^{2x}}{2} \, dx \\
    = -\frac{xe^{2x}}{4x+2} + \frac{e^{2x}}{4}+C = \boxed{\frac{e^{2x}}{8x+4} + C}
\end{gather*}

\problem{8}

\subproblem{a}
\begin{gather*}
    \int \sec x \, dx  = \int \frac{\sec^2 x+\sec x \tan x }{\sec x + \tan x }\, dx \\
    u=\sec x + \tan x\\
    du = \sec^2x +\sec x \tan x \, dx\\
    \int \frac{du}{u} = \ln |u|=\boxed{\ln |\sec x + \tan x | + C}
\end{gather*}

\subproblem{b}
\begin{gather*}
    \int \sec^3 x\, dx =
    \begin{bmatrix}
        \sec x & \sec^2x \, dx\\
        \sec x \tan x \, dx & \tan x
    \end{bmatrix}
    \sec x \tan x - \int \sec x \tan^2 x \, dx\\
    = \sec x \tan  x- \int \sec^3 x -\sec x \, dx\\
    2 \int \sec ^3 x\, dx= \sec x \tan x + \ln |\sec x + \tan x |\\
    \int \sec^3 x \,dx = \boxed{\frac{1}{2}\paren{\sec x \tan x + \ln | \sec x + \tan x |} +C}
\end{gather*}

\problem{9}

\subproblem{a}
\begin{gather*}
    J_n= \int _0^1 x^ne^{-x}\, dx = 
    \begin{bmatrix}
        x^n & e^{-x}\, dx\\
        nx^{n-1 }\, dx & -e^{-x}
    \end{bmatrix}
    -x^ne^{-x}\Big|^1_0 + n\int_0^1 x^{n-1}e^{-x}\,dx= \boxed{n J_{n-1} - \frac{1}{e}}\\
    J_0= \int_0^1  x^0 e^{-x}\, dx = \int_0^1 e^{-x}\, dx = -e^{-x} \Big|^1_0 = \boxed{1- \frac{1}{e}}
\end{gather*}

\subproblem{b}
\vem
\textit{Proof.}
\[
J_n=n!-\frac{n!}{e}\sum_{k=0}^n \frac{1}{k!},\quad n\ge 0
\]
Base case ($n=0$):
\[
J_0=0!- \frac{0!}{e} \sum_{k=0}^0 \frac{1}{k!} =1-\frac{1}{e}\cdot \frac{1}{1}= 1-\frac{1}{e}
\]
Inductive step: 

Assume the statement holds for $n=m,\quad m\ge 0$:
\[
J_m= m! - \frac{m!}{e}\sum_{k=0}^m \frac{1}{k!}
\]  
We now show that the statement holds for $n=m+1$:
\begin{align*}
    J_{m+1} &= (m+1)J_{m} - \frac{1}{e}\\
    &= (m+1)\paren{m! - \frac{m!}{e}\sum_{k=0}^m \frac{1}{k!}} - \frac{1}{e}\\
    &=(m+1)m!- \frac{(m+1)m!}{e} \sum_{k=0}^m \frac{1}{k!} - \frac{1}{e}\\
    &=(m+1)m!- \frac{(m+1)m!}{e} \paren{\sum_{k=0}^m \frac{1}{k!} + \frac{1}{(m+1)!}}\\
    &= (m+1)! - \frac{(m+1)!}{e} \sum_{k=0}^{m+1} \frac{1}{k!}
\end{align*}
Hence, the statement holds for $n=m+1$. Therefore, by the Principle of Mathematical Induction:
\[
J_n=n!-\frac{n!}{e}\sum_{k=0}^n \frac{1}{k!},\quad n\ge 0
\]

\subproblem{c}
\begin{gather*}
    J_n=\int_0^1 x^ne^{-x}\, dx\\
    x^n, e^{-x }>0 \implies J_n>0\\
    J_n=\int_0^1 x^ne^{-x}\, dx < \int_0^1 x^n\, dx = \frac{x^{n+1}}{n+1} \Bigg|^1_0=\frac{1}{n+1}\\
    \boxed{0 < J_n< \frac{1}{n+1}}
\end{gather*}

\subproblem{d}
\begin{gather*}
    let \, e=\frac{p}{q},n\ge \max (q,2)\\
    J_n=n!-\frac{n!}{e}\sum_{k=0}^n \frac{1}{k!}\\
    e\, J_n= e\cdot n! - n! \sum_{k=0}^n \frac{1}{k!}\\
    e\, J_n= \boxed{\frac{p}{q}\cdot n! - n! \sum_{k=0}^n \frac{1}{k!}}\\
    0 < J_n< \frac{1}{n+1}\\
    0 < e\, J_n< \frac{e}{n+1}\\
    n \ge 2 \implies n +1 \ge  3 \implies \frac{e}{n+1}<1\\
    \boxed{0< e\, J_n < 1}
\end{gather*}
As the value of $e\, J_n$ is within this range, it cannot be an integer, thus contradicting the assumption that $e$ is rational.\vem

\subproblem{e}
\begin{gather*}
    J_n=n!-\frac{n!}{e}\sum_{k=0}^n \frac{1}{k!}\\
    \sum_{k=0}^n \frac{1}{k!} = e - \frac{e \, J_n}{n!}\\
    0< e\, J_n < 1 \implies \lim_{n \to \infty }\frac{e \, J_n}{n!} =0\\
    \boxed{\lim_{n\to \infty} \sum_{k=0}^n \frac{1}{k!} =e}
\end{gather*}



















\end{document}

